\chapter{2009年大纲II卷（理）}

\section{单选题}

\begin{problem}
\(\frac{10\i}{2 - \i} =\)（\quad）

\choices
    {\(-2 + 4\i\)}
    {\(-2 - 4\i\)}
    {\(2 + 4\i\)}
    {\(2 - 4\i\)}
\answer{A}

\begin{solution}
分子、分母同乘以分母的共轭复数，得
\[
\frac{10\i}{2-\i}=\frac{10\i(2+\i)}{(2-\i)(2+\i)}
=\frac{20\i+10\i^{2}}{5}=-2+4\i
\]
故选 A.
\end{solution}
\end{problem}

\begin{problem}
设集合 \(A = \{x \mid x > 3\}\)，\(B = \left\{x \mid \frac{x - 1}{x - 4} < 0\right\}\)，则 \(A \cap B =\)（\quad）

\choices
    {\(\varnothing\)}
    {\((3,4)\)}
    {\((-2,1)\)}
    {\((4, + \infty)\)}
\answer{B}

\begin{solution}
由 \(\dfrac{x-1}{x-4}<0\)，临界点为 \(1\) 和 \(4\)；在 \((-\infty,1)\) 与 \((4,+\infty)\) 上分子、分母同号，商为正，在 \((1,4)\) 上分子、分母异号，商为负，且 \(x=4\) 不在定义域内，所以
\[
B=(1,4)
\]
因此 \(A\cap B=(3,+\infty)\cap(1,4)=(3,4)\)，故选 B.
\end{solution}
\end{problem}

\begin{problem}
已知 \(\triangle ABC\) 中，\(\cot A = -\frac{12}{5}\)，则 \(\cos A =\)（\quad）

\choices
    {\(\frac{12}{13}\)}
    {\(\frac{5}{13}\)}
    {\(-\frac{5}{13}\)}
    {\(-\frac{12}{13}\)}
\answer{D}

\begin{solution}
三角形的内角满足 \(0<A<\pi\)，而 \(\cot A<0\)，所以 \(A\) 为钝角，\(\cos A<0\)，\(\sin A>0\). 由
\[
\cot A=\frac{\cos A}{\sin A}=-\frac{12}{5}
\]
又 \(\sin^{2}A+\cos^{2}A=1\)，且 \(\sin A>0\)，故 \(\sin A=\frac{5}{13}\)，从而 \(\cos A=-\frac{12}{13}\). 故选 D.
\end{solution}
\end{problem}

\begin{problem}
曲线 \( y = \frac{x}{2x - 1} \) 在点 \((1,1)\) 处的切线方程为（\quad）

\choices
    {\(x - y - 2 = 0\)}
    {\(x + y - 2 = 0\)}
    {\(x + 4y - 5 = 0\)}
    {\(x - 4y - 5 = 0\)}
\answer{B}

\begin{solution}
设 \(f(x)=\dfrac{x}{2x-1}\)，则
\[
f'(x)=\frac{(2x-1)-2x}{(2x-1)^2}=-\frac{1}{(2x-1)^2}
\]
所以 \(f'(1)=-1\).
曲线在 \((1,1)\) 处的切线为 \(y-1=-(x-1)\)，即 \(x+y-2=0\). 故选 B.
\end{solution}
\end{problem}

\begin{problem}
已知正四棱柱 \(ABCD - A_{1}B_{1}C_{1}D_{1}\) 中，\(AA_{1} = 2AB\)，\(E\) 为 \(AA_{1}\) 中点，则异面直线 \(BE\) 与 \(CD_{1}\) 所成的角的余弦值为（\quad）

\choices
    {\(\frac{\sqrt{10}}{10}\)}
    {\(\frac{1}{5}\)}
    {\(\frac{3\sqrt{10}}{10}\)}
    {\(\frac{3}{5}\)}
\answer{C}

\begin{solution}
设 \(AB=1\)，则 \(AA_{1}=2\). 取坐标 \(A=(0,0,0)\)，\(B=(1,0,0)\)，\(C=(1,1,0)\)，\(D_{1}=(0,1,2)\)，\(E=(0,0,1)\).
于是 \(\overrightarrow{BE}=(-1,0,1)\)，可取 \(CD_{1}\) 的方向向量为 \(\overrightarrow{CD_{1}}=(-1,0,2)\). 两向量的夹角取锐角，故
\[
\cos\theta=\frac{\overrightarrow{BE}\cdot\overrightarrow{CD_{1}}}{|\overrightarrow{BE}|\,|\overrightarrow{CD_{1}}|}
=\frac{3}{\sqrt{2}\sqrt{5}}=\frac{3\sqrt{10}}{10}
\]
故选 C.
\end{solution}
\end{problem}

\begin{problem}
已知向量 \(\bs{a} = (2, 1)\)，\(\bs{a} \cdot \bs{b} = 10\)，\(|\bs{a} + \bs{b}| = 5\sqrt{2}\)，则 \(|\bs{b}| =\)（\quad）

\choices
    {\(\sqrt{5}\)}
    {\(\sqrt{10}\)}
    {\(5\)}
    {\(25\)}
\answer{C}

\begin{solution}
由 \(|\bs{a}|^{2}=2^{2}+1^{2}=5\)，以及
\[
|\bs{a}+\bs{b}|^{2}=|\bs{a}|^{2}+2\bs{a}\cdot\bs{b}+|\bs{b}|^{2}
\]
得 \(50=5+2\times10+|\bs{b}|^{2}\)，所以 \(|\bs{b}|^{2}=25\)，从而 \(|\bs{b}|=5\). 故选 C.
\end{solution}
\end{problem}

\begin{problem}
设 \(a = \log_3 \pi\)，\(b = \log_2 \sqrt{3}\)，\(c = \log_3 \sqrt{2}\)，则（\quad）

\choices
    {\(a > b > c\)}
    {\(a > c > b\)}
    {\(b > a > c\)}
    {\(b > c > a\)}
\answer{A}

\begin{solution}
因为 \(\pi>3\)，所以 \(a=\log_{3}\pi>1\). 又因为 \(\sqrt{3}<2\)，所以 \(b=\log_{2}\sqrt{3}<1\)，从而 \(a>b\). 此外
\[
b=\frac{\ln3}{2\ln2}
\]
又 \(c=\frac{\ln2}{2\ln3}\). 且 \(\ln3>\ln2>0\)，并且 \(b-c=\frac{(\ln3)^2-(\ln2)^2}{2\ln2\ln3}>0\)，故 \(b>c\). 因此 \(a>b>c\)，故选 A.
\end{solution}
\end{problem}

\begin{problem}
若将函数 \( y = \tan \left( \omega x + \frac{\pi}{4} \right) \)（\(\omega > 0\)）的图象向右平移 \( \frac{\pi}{6}\text{ 个单位长度} \)后，与函数 \( y = \tan \left( \omega x + \frac{\pi}{6} \right) \) 的图象重合，则 \( \omega \) 的最小值为（\quad）

\choices
    {\(\frac{1}{6}\)}
    {\(\frac{1}{4}\)}
    {\( \frac{1}{3} \)}
    {\(\frac{1}{2}\)}
\answer{D}

\begin{solution}
向右平移 \(\frac{\pi}{6}\) 后，函数变为
\[
y=\tan\left(\omega x+\frac{\pi}{4}-\frac{\omega\pi}{6}\right)
\]
两正切函数图象重合，当且仅当其自变量相差整数倍的 \(\pi\)，所以存在 \(k\in\mathbb Z\) 使
\[
\frac{\pi}{4}-\frac{\omega\pi}{6}=\frac{\pi}{6}+k\pi
\]
解得 \(\omega=\frac12-6k\).
由 \(\omega>0\)，最小值在 \(k=0\) 时取得，为 \(\frac12\). 故选 D.
\end{solution}
\end{problem}

\begin{problem}
已知直线 \( y = k(x + 2) \) （\( k > 0 \)） 与抛物线 \( C: y^2 = 8x \) 相交于 \(A\)，\(B\) 两点，\(F\) 为 \(C\) 的焦点. 若 \( |FA| = 2|FB| \)，则 \( k = \)（\quad）

\choices
    {\(\frac{1}{3}\)}
    {\(\frac{\sqrt{2}}{3}\)}
    {\( \frac{2}{3} \)}
    {\(\frac{2\sqrt{2}}{3}\)}
\answer{D}

\begin{solution}
抛物线 \(y^{2}=8x\) 的焦点为 \(F=(2,0)\)，准线为 \(x=-2\). 设 \(A\)，\(B\) 的横坐标分别为 \(x_{A}\)，\(x_{B}\). 由抛物线的定义，
\(|FA|=x_{A}+2\)，\(|FB|=x_{B}+2\).
将直线 \(y=k(x+2)\) 代入抛物线，得
\[
k^{2}x^{2}+(4k^{2}-8)x+4k^{2}=0
\]
故 \(x_{A}x_{B}=4\). 由 \(|FA|=2|FB|\) 得 \(x_{A}=2x_{B}+2\)，所以
\[
(2x_{B}+2)x_{B}=4
\]
因交点在抛物线上，\(x_{B}\geqslant0\)，故 \(x_{B}=1\)，\(x_{A}=4\). 再由根的和
\[
x_{A}+x_{B}=\frac{8}{k^{2}}-4
\]
得 \(5=\frac{8}{k^{2}}-4\)，于是 \(k=\frac{2\sqrt2}{3}\). 故选 D.
\end{solution}
\end{problem}

\begin{problem}
甲，乙两人从 \(4\text{ 门}\)课程中各选修 \(2\text{ 门}\)，则甲，乙所选的课程中至少有 \(1\text{ 门}\)不相同的选法共有（\quad）

\choices
    {\(6\text{ 种}\)}
    {\(12\text{ 种}\)}
    {\(30\text{ 种}\)}
    {\(36\text{ 种}\)}
\answer{C}

\begin{solution}
甲、乙两人各有 \(\mathrm C_4^2=6\) 种选法，因此全部选法有 \(6\times6=36\) 种. 两人所选课程完全相同的选法有 \(6\) 种. “至少有 \(1\) 门不相同”是排除完全相同的情形，故选法数为 \(36-6=30\)，故选 C.
\end{solution}
\end{problem}

\begin{problem}
已知双曲线 \(C: \frac{x^2}{a^2} - \frac{y^2}{b^2} = 1 \) （\(a > 0\)，\(b > 0\)）的右焦点为 \(F\)，过 \(F\) 且斜率为 \(\sqrt{3}\) 的直线交 \(C\) 于 \(A\)，\(B\) 两点，若 \(\overrightarrow{AF} = 4\overrightarrow{FB}\)，则 \(C\) 的离心率为（\quad）

\choices
    {\( \frac{6}{5} \)}
    {\( \frac{7}{5} \)}
    {\( \frac{5}{8} \)}
    {\( \frac{9}{5} \)}
\answer{A}

\begin{solution}
设 \(F=(c,0)\)，过焦点的直线用参数表示为
\[
(x,y)=(c+t,\sqrt3t)
\]
代入双曲线并整理，交点对应的参数是方程
\[
(b^{2}-3a^{2})t^{2}+2b^{2}ct+b^{4}=0
\]
的两根. 由 \(\overrightarrow{AF}=4\overrightarrow{FB}\)，可设两根为 \(-4t\) 和 \(t\). 由韦达定理，
\(-3t=\frac{2b^{2}c}{3a^{2}-b^{2}}\)，\(-4t^{2}=-\frac{b^{4}}{3a^{2}-b^{2}}\).
将第一式平方后与第二式联立，得
\[
16c^{2}=9(3a^{2}-b^{2})
\]
又 \(c^{2}=a^{2}+b^{2}\)，所以 \(25b^{2}=11a^{2}\)，从而
\[
e^{2}=\frac{c^{2}}{a^{2}}=1+\frac{b^{2}}{a^{2}}=1+\frac{11}{25}=\frac{36}{25}
\]
因 \(e>0\)，得 \(e=\frac65\)，故选 A.
\end{solution}
\end{problem}

\begin{problem}
纸制的正方体的六个面根据其方位分别标记为上，下，东，南，西，北. 现有沿该正方体的一些棱将正方体剪开，外面朝上展平，得到如图的平面图形，则标“\(\triangle\)”的面的方位是（\quad）

\FigureLayoutDeclare{img/2009/syllabus_paper_2_science/q12_fig1.png}{3.4cm}{33bp 36bp 34bp 36bp}
\bitmapfigure{img/2009/syllabus_paper_2_science/q12_fig1.png}

\choices
    {南}
    {北}
    {西}
    {下}
\answer{B}

\begin{solution}
将横向四个正方形依次记为 \(F_{0}\)，\(F_{1}\)，\(F_{2}\)，\(F_{3}\)，其中 \(F_{2}\) 标为“上”，\(F_{3}\) 标为“东”，而带三角形的正方形是 \(F_{0}\) 上方的面. 折叠后，横向四面依次围成正方体的一周，因此 \(F_{0}\) 与 \(F_{2}\) 相对，\(F_{1}\) 与 \(F_{3}\) 相对. 所以 \(F_{0}\) 的方位为“下”，\(F_{1}\) 的方位为“西”.

以 \(F_{2}\) 为“上”、右侧相邻的 \(F_{3}\) 为“东”折叠，\(F_{0}\) 上方的面折起后位于“北”侧；其下方的面位于“南”侧. 故带三角形的面方位为“北”，选项为 B.
\end{solution}
\end{problem}

\section{填空题}

\begin{problem}
\(\left(x\sqrt{y} - y\sqrt{x}\right)^4\) 的展开式中 \(x^3 y^3\) 的系数为\fillinblank{}.
\answer{6}

\begin{solution}
展开式通项中，若从四个因式中取出 \(k\) 个 \(-y\sqrt{x}\)，则该项为
\[
(-1)^k\mathrm C_4^k x^{4-k/2}y^{2+k/2}
\]
令 \(4-\frac{k}{2}=3\)，得 \(k=2\)，此时 \(y\) 的指数也为 \(2+\frac{k}{2}=3\). 因而所求系数为 \(\mathrm C_4^2=6\).
\end{solution}
\end{problem}

\begin{problem}
设等差数列 \(\{a_{n}\}\) 的前 \(n\) 项和为 \(S_{n}\)，若 \(a_{5} = 5a_{3}\)，则 \(\frac{S_{9}}{S_{5}} = \)\fillinblank{}.
\answer{9}

\begin{solution}
设首项为 \(a_{1}\)，公差为 \(d\). 由 \(a_{5}=5a_{3}\)，得 \(a_{1}+4d=5(a_{1}+2d)\)，即 \(2a_{1}+3d=0\).
于是 \(a_{1}=-\frac32d\)，并且
\[
S_{9}=\frac92(2a_{1}+8d)=\frac{45}{2}d
\]
同理，\(S_{5}=\frac52(2a_{1}+4d)=\frac{5}{2}d\).
题目中的比值有意义，故 \(d\ne0\)，从而 \(\dfrac{S_{9}}{S_{5}}=9\).
\end{solution}
\end{problem}

\begin{problem}
设 \(OA\) 是球 \(O\) 的半径，\(M\) 是 \(OA\) 的中点，过 \(M\) 且与 \(OA\) 成 \(45^{\circ}\) 角的平面截球 \(O\) 的表面得到圆 \(C\). 若圆 \(C\) 的面积等于 \(\frac{7\pi}{4}\)，则球 \(O\) 的表面积等于\fillinblank{}.
\answer{\(8\pi\)}

\begin{solution}
设球的半径为 \(R\)，球心到截面平面的距离为 \(d\)，截面圆半径为 \(r\). 因为 \(OM=\frac{R}{2}\)，且直线 \(OA\) 与平面的夹角为 \(45^{\circ}\)，所以
\[
d=OM\sin45^{\circ}=\frac{R}{2\sqrt2}
\]
由截面圆面积得 \(\pi r^{2}=\frac{7\pi}{4}\)，即 \(r^{2}=\frac74\). 在球心、截面圆心和截面圆上任一点构成的直角三角形中，
\[
R^{2}=d^{2}+r^{2}=\frac{R^{2}}8+\frac74
\]
故 \(R^{2}=2\). 所以球的表面积为 \(4\pi R^{2}=8\pi\).
\end{solution}
\end{problem}

\begin{problem}
已知 \(AC\)，\(BD\) 为圆 \(O: x^{2} + y^{2} = 4\) 的两条相互垂直的弦，垂足为 \(M(1, \sqrt{2})\)，则四边形 \(ABCD\) 的面积的最大值为\fillinblank{}.
\answer{5}

\begin{solution}
设圆心到弦 \(AC\)，\(BD\) 所在直线的距离分别为 \(d_{1}\)，\(d_{2}\). 因两直线垂直且交于 \(M\)，有
\[
d_{1}^{2}+d_{2}^{2}=OM^{2}=1^{2}+\left(\sqrt2\right)^{2}=3
\]
圆半径为 \(2\)，故两弦长分别为 \(2\sqrt{4-d_{1}^{2}}\) 和 \(2\sqrt{4-d_{2}^{2}}\). 四边形 \(ABCD\) 的对角线为这两条互相垂直的弦，所以其面积
\[
S=\frac12 AC\cdot BD=2\sqrt{(4-d_{1}^{2})(4-d_{2}^{2})}
\]
令 \(u=d_{1}^{2}\)，\(v=d_{2}^{2}\)，则 \(u+v=3\)，且
\[
(4-u)(4-v)=4+uv\leqslant4+\left(\frac{u+v}{2}\right)^{2}=\frac{25}{4}
\]
因此 \(S\leqslant2\sqrt{\frac{25}{4}}=5\). 当且仅当 \(u=v=\frac32\) 时取等号；取过 \(M\) 且分别与 \(OM\) 成 \(45^{\circ}\) 角的两条互相垂直直线即可实现，故最大值为 \(5\).
\end{solution}
\end{problem}

\section{解答题}

\begin{problem}
设 \(\triangle ABC\) 的内角 \(A\)，\(B\)，\(C\) 的对边长分别为 \(a\)，\(b\)，\(c\)，\(\cos (A - C) + \cos B = \frac{3}{2}\)，\(b^2 = ac\)，求 \(B\).

\begin{solution}
由 \(A+B+C=\pi\)，得 \(B=\pi-(A+C)\)，所以 \(\cos B=-\cos(A+C)\). 因而
\[
\cos(A-C)-\cos(A+C)=\frac32
\]
利用和差化积公式，得 \(2\sin A\sin C=\frac32\)，即
\[
\sin A\sin C=\frac34
\]
由正弦定理，存在 \(R>0\) 使 \(a=2R\sin A\)，\(b=2R\sin B\)，\(c=2R\sin C\). 代入 \(b^{2}=ac\)，得
\[
\sin^{2}B=\sin A\sin C=\frac34
\]
因为 \(0<B<\pi\)，所以 \(\sin B=\frac{\sqrt3}{2}\)，从而 \(B=\frac\pi3\) 或 \(B=\frac{2\pi}{3}\).

又由 \(b^{2}=ac\)，至少有一个数 \(a\)，\(c\) 不小于 \(b\). 由边角大小规律，至少有一个角 \(A\)，\(C\) 不小于 \(B\)，于是 \(A+C\geqslant B\)，即 \(\pi-B\geqslant B\)，所以 \(B\leqslant\frac\pi2\). 排除 \(B=\frac{2\pi}{3}\)，故
\[
B=\frac\pi3=60^{\circ}
\]
\end{solution}
\end{problem}

\begin{problem}
如图，直三棱柱 \(ABC - A_{1}B_{1}C_{1}\) 中，\(AB \perp AC\)，\(D\)，\(E\) 分别为 \(AA_{1}\)，\(B_{1}C\) 的中点，\(DE \perp\) 平面 \(BCC_{1}\).

\begin{enumerate}
\item 证明：\(AB = AC\).
\item 设二面角 \(A - BD - C\) 为 \(60^{\circ}\)，求 \(B_{1}C\) 与平面 \(BCD\) 所成的角的大小.
\end{enumerate}

\bitmapfigure[width=0.36\linewidth]{img/2009/syllabus_paper_2_science/q18_fig1.png}

\begin{solution}
\begin{enumerate}
\item 设 \(AB=u\)，\(AC=v\)，\(AA_{1}=h\). 建立空间直角坐标系，使
\(A=(0,0,0)\)，\(B=(u,0,0)\)，\(C=(0,v,0)\)，\(A_{1}=(0,0,h)\).
则 \(D=(0,0,\frac h2)\)，\(B_{1}=(u,0,h)\)，且
\[
\overrightarrow{DE}=\left(\frac u2,\frac v2,0\right)
\]
平面 \(BCC_{1}\) 的一个法向量为 \((v,u,0)\). 由 \(DE\perp\) 平面 \(BCC_{1}\)，有 \((u,v,0)\parallel(v,u,0)\)，所以 \(u=v\)，即 \(AB=AC\).
\item 记 \(AB=AC=s\). 此时平面 \(ABD\) 的一个法向量可取 \(\bs{n}_{1}=(0,1,0)\)，平面 \(BCD\) 的一个法向量可取
\[
\bs{n}_{2}=\left(\frac h2,\frac h2,s\right)
\]
二面角为 \(60^{\circ}\)，故
\[
\frac{\left|\bs{n}_{1}\cdot\bs{n}_{2}\right|}{|\bs{n}_{1}|\,|\bs{n}_{2}|}
=\frac{h/2}{\sqrt{h^{2}/2+s^{2}}}=\cos60^{\circ}=\frac12
\]
从而 \(h^{2}=2s^{2}\).

直线 \(B_{1}C\) 的方向向量可取 \(\overrightarrow{B_{1}C}=(-s,s,-h)\). 设它与平面 \(BCD\) 所成的角为 \(\theta\)，则
\[
\sin\theta=\frac{\left|\overrightarrow{B_{1}C}\cdot\bs{n}_{2}\right|}{|\overrightarrow{B_{1}C}|\,|\bs{n}_{2}|}
=\frac{hs}{\sqrt{2s^{2}+h^{2}}\sqrt{h^{2}/2+s^{2}}}
=\frac12
\]
线面角取锐角，故 \(\theta=30^{\circ}\).
\end{enumerate}
\end{solution}
\end{problem}

\begin{problem}
设数列 \(\{a_{n}\}\) 的前 \(n\) 项和为 \(S_{n}\)，已知 \(a_{1} = 1\)，\(S_{n+1} = 4a_{n} + 2\).

\begin{enumerate}
\item 设 \(b_{n} = a_{n+1} - 2a_{n}\). 证明数列 \(\{b_n\}\) 是等比数列.
\item 求数列 \(\{a_{n}\}\) 的通项公式.
\end{enumerate}

\begin{solution}
\begin{enumerate}
\item 由已知条件在 \(n=1\) 时给出
\[
a_{1}+a_{2}=S_{2}=4a_{1}+2=6
\]
所以 \(a_{2}=5\)，从而 \(b_{1}=a_{2}-2a_{1}=3\). 对任意 \(n\geqslant1\)，相邻两式相减得
\[
a_{n+2}=S_{n+2}-S_{n+1}=4a_{n+1}+2-(4a_{n}+2)=4(a_{n+1}-a_{n})
\]
因此
\[
b_{n+1}=a_{n+2}-2a_{n+1}=2(a_{n+1}-2a_{n})=2b_{n}
\]
所以 \(\{b_{n}\}\) 是首项为 \(3\)、公比为 \(2\) 的等比数列，即 \(b_{n}=3\cdot2^{n-1}\).
\item 由 \(b_{n}=a_{n+1}-2a_{n}\)，有
\[
\frac{a_{n+1}}{2^{n}}-\frac{a_{n}}{2^{n-1}}=\frac{b_{n}}{2^{n}}=\frac32
\]
数列 \(\left\{\frac{a_{n}}{2^{n-1}}\right\}\) 是首项为 \(1\)、公差为 \(\frac32\) 的等差数列，故
\[
\frac{a_{n}}{2^{n-1}}=1+(n-1)\frac32=\frac{3n-1}{2}
\]
所以数列 \(\{a_n\}\) 的通项公式为
\(a_n=(3n-1)2^{n-2}\)，\((n\geqslant1)\).
\end{enumerate}
\end{solution}
\end{problem}

\begin{problem}
某车间甲组有 \(10\text{ 名}\)工人，其中有 \(4\text{ 名}\)女工人；乙组有 \(5\text{ 名}\)工人，其中有 \(3\text{ 名}\)女工人. 现采用分层抽样方法（层内采用不放回简单随机抽样）从甲，乙两组中共抽取 \(3\text{ 名}\)工人进行技术考核.

\begin{enumerate}
\item 求从甲，乙两组各抽取的人数.
\item 求从甲组抽取的工人中恰有 \(1\text{ 名}\)女工人的概率.
\item 记 \(\xi\) 表示抽取的 \(3\text{ 名}\)工人中男工人数，求 \(\xi\) 的分布列及数学期望.
\end{enumerate}

\begin{solution}
\begin{enumerate}
\item 甲、乙两组人数之比为 \(10:5=2:1\). 按分层抽样的比例，从两组共抽取 \(3\) 人，所以甲组抽取 \(2\) 人，乙组抽取 \(1\) 人.
\item 甲组有 \(4\) 名女工人和 \(6\) 名男工人，故所求概率为
\[
\frac{\mathrm C_4^1\mathrm C_6^1}{\mathrm C_{10}^{2}}=\frac{24}{45}=\frac{8}{15}
\]
\item 设 \(X\) 为从甲组抽取的男工人数，\(Y\) 为从乙组抽取的男工人数，则 \(\xi=X+Y\). 由不放回抽样，
\[
P(X=0)=\frac{\mathrm C_4^2}{\mathrm C_{10}^{2}}=\frac{2}{15}
\]
\[
P(X=1)=\frac{\mathrm C_6^1\mathrm C_4^1}{\mathrm C_{10}^{2}}=\frac{8}{15}
\]
\[
P(X=2)=\frac{\mathrm C_6^2}{\mathrm C_{10}^{2}}=\frac13
\]
并且 \(P(Y=0)=\frac35\)，\(P(Y=1)=\frac25\). 两组抽样相互独立，因此
\[
P(\xi=0)=\frac{2}{15}\cdot\frac35=\frac{6}{75}
\]
\[
P(\xi=1)=\frac{2}{15}\cdot\frac25+\frac{8}{15}\cdot\frac35=\frac{28}{75}
\]
\[
P(\xi=2)=\frac{8}{15}\cdot\frac25+\frac13\cdot\frac35=\frac{31}{75}
\]
\[
P(\xi=3)=\frac13\cdot\frac25=\frac{10}{75}
\]
所以分布列为
\begin{center}
\begin{tabular}{c|cccc}
\(\xi\) & 0 & 1 & 2 & 3 \\
\hline
\(P\) & \(\frac{6}{75}\) & \(\frac{28}{75}\) & \(\frac{31}{75}\) & \(\frac{10}{75}\)
\end{tabular}
\end{center}
数学期望为
\[
E(\xi)=0\cdot\frac{6}{75}+1\cdot\frac{28}{75}+2\cdot\frac{31}{75}+3\cdot\frac{10}{75}=\frac85
\]
\end{enumerate}
\end{solution}
\end{problem}

\begin{problem}
已知椭圆 \(C: \frac{x^2}{a^2} + \frac{y^2}{b^2} = 1 \) （\(a > b > 0\)）的离心率为 \(\frac{\sqrt{3}}{3}\)，过右焦点 \(F\) 的直线 \(l\) 与 \(C\) 相交于 \(A\)，\(B\) 两点，当 \(l\) 的斜率为 \(1\) 时，坐标原点 \(O\) 到 \(l\) 的距离为 \(\frac{\sqrt{2}}{2}\).

\begin{enumerate}
\item 求 \(a\)，\(b\) 的值.
\item \(C\) 上是否存在点 \(P\)，使得当 \(l\) 绕 \(F\) 转到某一位置时，有 \(\overrightarrow{OP} = \overrightarrow{OA} + \overrightarrow{OB}\) 成立？若存在，求出所有的 \(P\) 的坐标与 \(l\) 的方程；若不存在，说明理由.
\end{enumerate}

\begin{solution}
\begin{enumerate}
\item 设椭圆的半焦距为 \(c\)，则右焦点为 \(F=(c,0)\)，且
\[
\frac{c}{a}=\frac{\sqrt3}{3}
\]
又 \(b^{2}=a^{2}-c^{2}\).
当直线 \(l\) 的斜率为 \(1\) 时，其方程为 \(y=x-c\)，即 \(x-y-c=0\). 由原点到直线的距离，
\[
\frac{c}{\sqrt2}=\frac{\sqrt2}{2}
\]
得 \(c=1\). 因而 \(a=\sqrt3\)，且 \(b^{2}=3-1=2\)，所以 \(b=\sqrt2\).
\item 由上知椭圆方程为 \(2x^{2}+3y^{2}=6\)，焦点为 \(F=(1,0)\). 先设 \(l\) 不是竖直直线，写成
\[
y=k(x-1)
\]
设它与椭圆交于 \(A(x_{1},y_{1})\)，\(B(x_{2},y_{2})\). 将直线方程代入椭圆，得
\[
(2+3k^{2})x^{2}-6k^{2}x+3k^{2}-6=0
\]
由韦达定理，记 \(D=2+3k^{2}\)，则
\[
x_{1}+x_{2}=\frac{6k^{2}}{D}
\]
又 \(y_{1}+y_{2}=k(x_{1}+x_{2}-2)=-\frac{4k}{D}\).
若 \(\overrightarrow{OP}=\overrightarrow{OA}+\overrightarrow{OB}\)，则
\[
P=\left(\frac{6k^{2}}{D},-\frac{4k}{D}\right)
\]
点 \(P\) 在椭圆上等价于
\[
2\left(\frac{6k^{2}}{D}\right)^{2}+3\left(-\frac{4k}{D}\right)^{2}=6
\]
即 \(\frac{24k^{2}}{D}=6\). 因而 \(4k^{2}=2+3k^{2}\)，所以 \(k^{2}=2\).

当 \(k=\sqrt2\) 时，\(P=\left(\frac32,-\frac{\sqrt2}{2}\right)\)，\(l:y=\sqrt2(x-1)\)；当 \(k=-\sqrt2\) 时，\(P=\left(\frac32,\frac{\sqrt2}{2}\right)\)，\(l:y=-\sqrt2(x-1)\).

最后检查竖直直线 \(x=1\). 它与椭圆的两个交点纵坐标互为相反数，故两点坐标之和为 \((2,0)\)，而 \(2(2)^{2}+3(0)^{2}=8\ne6\)，所以 \(P=(2,0)\) 不在椭圆上. 因此以上两点即为全部满足条件的点.
\end{enumerate}
\end{solution}
\end{problem}

\begin{problem}
设函数 \( f(x) = x^{2} + a\ln (1 + x) \) 有两个极值点 \(x_{1}\)，\(x_{2}\)，且 \( x_{1} < x_{2} \).

\begin{enumerate}
\item 求 \(a\) 的取值范围，并讨论 \(f(x)\) 的单调性.
\item 证明： \(f(x_{2}) > \frac{1 - 2\ln 2}{4}\).
\end{enumerate}

\begin{solution}
\begin{enumerate}
\item 函数的定义域为 \((-1,+\infty)\)，且
\[
f'(x)=2x+\frac{a}{1+x}=\frac{2x^{2}+2x+a}{1+x}
\]
要有两个定义域内的极值点，二次方程 \(2x^{2}+2x+a=0\) 必须有两个不相等且都大于 \(-1\) 的根. 其根为
\[
x_{1}=\frac{-1-\sqrt{1-2a}}{2}
\]
又 \(x_{2}=\frac{-1+\sqrt{1-2a}}{2}\).
两根不相等要求 \(a<\frac12\)，而 \(x_{1}>-1\) 要求 \(\sqrt{1-2a}<1\)，即 \(a>0\). 因此
\[
0<a<\frac12
\]
在定义域内 \(1+x>0\)，而二次式 \(2x^{2}+2x+a\) 在两根外为正、两根之间为负，所以 \(f(x)\) 在 \((-1,x_{1})\) 和 \((x_{2},+\infty)\) 上单调递增，在 \((x_{1},x_{2})\) 上单调递减.
\item 令 \(t=1+x_{2}\). 由 \(x_{2}\in(-1,0)\) 且 \(0<a<\frac12\)，有 \(\frac12<t<1\). 又因 \(f'(x_{2})=0\)，
\[
a=-2x_{2}(1+x_{2})=2t(1-t)
\]
于是
\[
f(x_{2})=(1-t)^{2}+2t(1-t)\ln t=:H(t)
\]
求导得
\[
H'(t)=2(1-2t)\ln t
\]
当 \(\frac12<t<1\) 时，\(1-2t<0\) 且 \(\ln t<0\)，故 \(H'(t)>0\). 因而
\[
f(x_{2})=H(t)>H\left(\frac12\right)=\frac14+\frac12\ln\frac12=\frac{1-2\ln2}{4}
\]
\end{enumerate}
\end{solution}
\end{problem}
